\section{An effective obstruction for every fixed ordinary alternating tail}
\label{sec:effective-fixed-tail}

\paragraph{Status, origin and exact scope.}
New derivation D-AFT01--03, Q2ct. Complete author and root
whole-proof reviews and the separate /root/imprimitive proof,
primary-source and independent coefficient reconstruction have passed.
The cyclic character-defect method originates in the root's
HD-AT301 proof. The exact tensor-cycle and actual quotient-cover
framework is D-ATT07--09. The new assertion here is uniform over
every fixed tail partition, with an explicit computable threshold;
it does not assert a threshold uniform over all tail sizes.

Fix a partition $\mu\vdash k$, where $k\ge3$. For sufficiently
large $n$, let $W_n=S^{(n-k,\mu)}$ be the rational ordinary
$S_n$ representation. We consider an actual rank-nine sixfold
inducing source whose irreducibly acting ordinary core is $A_n^t$
on $W_n^{\otimes t}$, with the inducing image transitive on the
simple factors. All compatible scalar and outer actions and all
factor permutations are allowed. For $t=1$ the conclusion holds
over every original imaginary quadratic field. For $t\ge2$ we
assume that the inducing character field contains a non-Eisenstein
imaginary quadratic field. Equality with the original field is
not assumed.

\subsection*{D-AFT01. The coefficient constants and uniform spectrum bound}

The published character-polynomial formula of Garsia--Goupil,
Proposition~I.1, printed p.~2, with its complete proof in
equations (2.1)--(2.6), pp.~9--11, gives
\[
 \chi_{W_n}(x)=q_\mu(X_1(x),\ldots,X_k(x))
 \quad(n-k\ge\mu_1).                                  \tag{AFT.1}
\]
Its exact rational polynomial is computed by the displayed
finite umbral formula in that proposition, or its recursive
Corollary~I.1. We use this published formula, not an unverified
asymptotic character assertion \cite{HD:GarsiaGoupil2009}.

Write $f=X_1$ and expand the polynomial in ordinary monomials:
\[
 q_\mu=Cf^k+R(f,X_2,\ldots,X_k),\qquad
 C=\frac{f^\mu}{k!}>0.                                 \tag{AFT.2}
\]
Every monomial has weighted degree at most $k$, where $X_i$
has weight $i$. The displayed leading term is the only monomial
of ordinary total degree $k$. Indeed each product
$\prod_i(iX_i-1)^{a_i}$ in the published formula has weighted
degree at most $\sum ia_i=k$, and the falling-factorial
substitution only lowers this degree. Ordinary degree $k$
therefore requires the monomial $X_1^k$, arising from the
identity class of $S_k$ with coefficient $f^\mu/k!$.

Let $r_\alpha$ be the coefficients of $R$, including its
constant coefficient, and set
\[
 \begin{aligned}
 M&=1+C^{-1}\sum_\alpha|r_\alpha|\ge1,
 &F&=\lceil kM\rceil+1,\\
 N_\mu&=\left\lceil\max\left\{
 2k+2,\ 200,\ \frac{100}{C},\ 100kM,
 [100(1+M)]^5,\ [100(1+M)F^k]^2
 \right\}\right\rceil .
 \end{aligned}                                         \tag{AFT.3}
\]
These are explicit rational computations and integer ceilings
from a finite published polynomial. In particular the threshold
does not conceal an ineffective asymptotic existence statement.

For every integer $n\ge N_\mu$, define
\[
 \begin{gathered}
 a=q_\mu(n,0,\ldots,0)=\dim W_n,\qquad S=Cn^{k-1},\\
 b=S/100,\qquad c=a/n-Cn^{k-2}/100,\\
 A=a/n-101Cn^{k-2}/100,\qquad B=2S/5,
 \qquad d_0=A+B.
 \end{gathered}                                        \tag{AFT.4}
\]
The symbol $S$ in (AFT.4) is a positive real number, not a
finite group. All displayed bounds below are rational for
integer $n$.

\begin{theorem}[D-AFT01]
For every nonidentity permutation $x\in S_n$, with natural
index $I(x)=n-\#\mathrm{cycles}(x)$,
\[
 \nu(W_nx)\ge AI(x)+B.                                 \tag{AFT.5}
\]
Here $\nu$ is maximum-eigenspace codimension, so the estimate
holds for every target eigenvalue and survives every scalar
twist. For every permutation, including the identity,
\[
 \chi_{W_n}(x)\ge-b.                                   \tag{AFT.6}
\]
Moreover
\[
 \begin{gathered}
 S\ge100,\qquad \tfrac{99}{100}nS\le a\le\tfrac{101}{100}nS,
 \qquad \tfrac9{10}S\le A\le\tfrac{101}{100}S,\\
 \tfrac{13}{10}S\le d_0\le\tfrac{141}{100}S,
 \qquad B\ge40,\quad b\le d_0,\quad 2d_0\le a-1.
 \end{gathered}                                        \tag{AFT.7}
\]
\end{theorem}

\begin{proof}
Write $q_\mu=Cf^k+P(f)+Q(f,X_2,\ldots,X_k)$, where $P$
contains the lower pure-$f$ terms and each monomial of $Q$
contains at least one other variable. Their total absolute
coefficient mass is at most $CM$.
Since $|P(n)|\le CMn^{k-1}$ and $M/n\le1/100$, the degree
bounds in (AFT.7) hold. The threshold also gives $S\ge100$.

First prove the global negative-character bound. If
$0\le f\le n^{3/5}$, a monomial $f^r\prod_{i\ge2}X_i^{a_i}$,
with $s=\sum_{i\ge2}a_i$, satisfies $r+2s\le k$ and has
absolute value at most $n^{3r/5+s}\le n^{3k/5}$.
The same estimate holds for the leading term. Therefore
\[
 |q_\mu|\le C(1+M)n^{3k/5}\le S/100.                  \tag{AFT.8}
\]
The last inequality uses $k\ge3$ and
$n^{1/5}\ge100(1+M)$.
If $f\ge n^{3/5}$, each lower pure term divided by $f^k$
is at most $n^{-3/5}$, and each mixed monomial divided by
$f^k$ is at most
$(n/f^2)^s\le n^{-1/5}$. Hence
$|R|\le CMn^{-1/5}f^k<Cf^k/100$ and $q_\mu>0$.
This proves (AFT.6) in all cases.

We next prove the particularly sharp global defect estimate
\[
             a-\chi_{W_n}(x)\ge c(n-f).                \tag{AFT.9}
\]
Put $m=n-f$. Every mixed monomial has ordinary degree at
most $k-1$. Selecting one non-$f$ variable and using
$X_i\le m/i\le m$, while bounding the remaining variables
by $n$, gives absolute value at most $mn^{k-2}$.
For a pure lower monomial of degree $r\ge2$,
\[
 \left|\frac fn n^r-f^r\right|
 =f(n^{r-1}-f^{r-1})\le(k-2)mn^{k-2}.
\]
For $r=1$ the difference is zero; for $r=0$ it is $m/n$,
also bounded by $(k-2)mn^{k-2}$. Consequently
\[
 \frac fn a-q_\mu
 \ge Cfm n^{k-2}-CM(k-2)mn^{k-2}\ge0
 \quad\hbox{if }f\ge F.                               \tag{AFT.10}
\]
Here the first positive term follows from
$fn^{k-1}-f^k\ge fm n^{k-2}$. The identity $m=0$ is
included without division.

For the remaining integer values $0\le f<F$, the weighted
degree bound instead gives
\[
 |q_\mu|\le C(1+M)F^k n^{k/2}\le S/100,                \tag{AFT.11}
\]
because $\sqrt n\ge100(1+M)F^k$ and
$(k+1)/2\le k-1$. If $f\ge1$, the lower degree bound
gives $af/n\ge99S/100$, so again $q_\mu\le af/n$.
If $f=0$, (AFT.11) gives $q_\mu\le S/100$.
These cases prove (AFT.9), with exactly the $c$ in (AFT.4).
The use of integer $f\ge1$ here is explicit.

We also need a stronger defect when many letters are fixed.
For a pure term of degree $r\le k-1$,
$|n^r-f^r|\le(k-1)mn^{k-2}$; the constant cancels.
Together with the mixed-monomial estimate already proved,
this gives
\[
 a-q_\mu\ge (n-f)c_f,\qquad
 c_f=C\left(\sum_{j=0}^{k-1}n^{k-1-j}f^j
                  -(k-1)Mn^{k-2}\right).               \tag{AFT.12}
\]
The function $c_f$ is nondecreasing for $f\ge0$. The upper
degree estimate before replacing $M/n$ by $1/100$ yields
\[
                 c_f-A\ge Cn^{k-2}(f-kM).              \tag{AFT.13}
\]
Indeed retain just the $j=0,1$ terms in the sum, use
$a/n\le Cn^{k-1}+CMn^{k-2}$, and discard the positive
$101Cn^{k-2}/100$.

Let $x$ have order $L$ and index $I\ge1$. If $m_z$ is a
target-eigenvalue multiplicity, cyclic averaging in (AFT.9)
gives $a-m_1\ge cI$. The real character bound (AFT.6) gives
\[
 m_z-m_1=\frac1L\sum_{j=0}^{L-1}\chi(x^j)
                    (\operatorname{Re}z^{-j}-1)\le b
 \quad(z\ne1);                                        \tag{AFT.14}
\]
the trivial target needs no estimate. Thus $\nu\ge cI-b$.
If $I\ge9n/20$, then $c-A=Cn^{k-2}$ implies
\[
                  \nu-AI\ge(9/20-1/100)S=11S/25>B.
                                                               \tag{AFT.15}
\]

Suppose instead $I<9n/20$. Every power of $x$ fixes at least
$f_0=n-2I>n/10$ letters, because the moved support is at most
twice the natural index. Averaging (AFT.12) gives
$a-m_1\ge c_{f_0}I$. Since $kM\le n/100$, (AFT.13)
gives $c_{f_0}-A\ge9S/100$. If $I\ge5$, this again yields
$\nu-AI\ge11S/25>B$. If $1\le I\le4$, then
$f_0\ge n-8\ge9n/10$ and the stronger bound is
$c_{f_0}-A\ge89S/100$. Therefore
$\nu-AI\ge88S/100>B$. This covers all nonidentity
permutations and every target eigenvalue.

Finally $n\ge200$ gives
$A\ge(99/100-101/(100n))S\ge9S/10$; its upper bound is
immediate. Adding $B$ proves the $d_0$ bounds. Since
$S\ge100$, $n\ge200$, and $a\ge99nS/100$, all remaining
comparisons in (AFT.7) follow directly. In particular
$a\ge19800$, $d_0\ge130$, and $2d_0\le a-1$.
\end{proof}

\subsection*{D-AFT02. One ordinary core in every field}

\begin{theorem}[D-AFT02]
For $n\ge N_\mu$, a rank-nine sixfold inducing source
cannot have the single ordinary core $A_n$ on $W_n$.
This holds over every original imaginary quadratic field,
with all compatible scalar and outer actions.
\end{theorem}

\begin{proof}
The threshold $n\ge2k+2$ makes $(n-k,\mu)$ non-self-conjugate:
its first row exceeds its number of rows. The restriction
to $A_n$ is therefore irreducible. It is nontrivial and
faithful because $A_n$ is simple. Its rational $S_n$ model
has trivial scalar intersection: a scalar centralizes the
faithful $A_n$ image, while $C_{S_n}(A_n)=1$.
The accepted ordinary alternating normalizer interface
therefore extracts an actual permutation homomorphism.
Its image contains $A_n$, so the associated degree-$n$
quotient cover is connected.

The exact intermediate Euler identity is
$6=2a(h-1)+D$. Since $a\ge19800$, genus at least two is
impossible. At genus one there are at most six genuine
marks and the actual period rank is at most six. Thus
rank nine requires $h=0$, at least twelve genuine inertia
entries, and $D=2a+6$. Riemann--Hurwitz for the connected
degree-$n$ quotient gives $\sum I(x_i)\ge2n-2$.

Every nonscalar inertia moves at least $AI(x_i)+B$
dimensions by D-AFT01. A genuine scalar has full moving
rank $a\ge B$ and index zero. Hence
\[
 \begin{split}
 D-2a&\ge2A(n-1)+12B-2a\\
     &=12B-\left(\frac{2a}{n}
             +\frac{202}{100}C(n-1)n^{k-2}\right)\\
     &\ge\left(\frac{24}{5}-\frac{101}{25}\right)S
       =\frac{19}{25}S\ge76>6.
 \end{split}                                           \tag{AFT.16}
\]
This contradicts the exact source budget. No condition on
the scalar cyclotomic field is needed for this case.
\end{proof}

\subsection*{D-AFT03. Every transitive tensor count}

\begin{theorem}[D-AFT03]
For $n\ge N_\mu$ and every $t\ge2$, a finite inducing
image normalizing $A_n^t$ on $W_n^{\otimes t}$ transitively
cannot give a rank-nine sixfold source if its character
field contains a non-Eisenstein imaginary quadratic field.
All compatible scalar phases, outer actions and factor
permutations are included.
\end{theorem}

\begin{proof}
Put $e=a^t$. The rational tensor-wreath normalizer meets
the scalars trivially by its conjugation action on the
faithful simple factors. Thus every inducing operator
has actual scalar and rational coordinates, and the
natural $nt$-letter quotient is transitive. The inducing
character field is contained in the cyclotomic field of
the scalar character. If this cyclotomic field contains
a non-Eisenstein imaginary quadratic field, its scalar
image is not contained in $\mu_6$. Generation and product
one in the scalar quotient consequently force at least
two inertia phases outside $\mu_6$. Only field inclusion
is used. The actual intermediate-cover Euler and Hurwitz
arguments give $h=0$, $r\ge12$, and $D=2e+6$.

Every power of a rational tensor-wreath operator has trace
at least $-(b/a)e$. Indeed its trace is a product of traces
of the factor-cycle products. A negative product contains
one negative factor of absolute value at most $b$, and
all other factors have absolute value at most $a$.
The rational cyclotomic multiplicity bound, together with
the fourth-root square-trace argument of D-ATT02, therefore
gives moving rank at least $(3-b/a)e/4$ at the two marked
scalar phases.

Every genuine inertia has moving rank at least $(d_0/a)e$.
For an operator fixing factors, use D-AFT01 on any active
factor and spectral convolution; a genuine scalar has
full rank. For an operator moving factors, the exact
weighted-word estimate gives the stronger
$(a-1)e/(2a)$, which is at least $(d_0/a)e$ by (AFT.7).
If at least three entries move factors, one is distinct
from both scalar marks. With nine other genuine entries,
\[
 D\ge2e+a^{t-1}\left(9d_0-\frac{b+1}{2}\right)>2e+6,
                                                               \tag{AFT.17}
\]
using $b\le d_0$ and $d_0\ge130$. Thus exactly two entries
move factors: zero contradicts transitivity and one
contradicts product one. Their inverse factor permutations
are full $t$-cycles.

Let $H$ be the total natural index of their two cycle
products, and $N$ the sum of the natural factor indices
at the remaining entries. Riemann--Hurwitz for the actual
connected degree-$nt$ quotient gives
$2n(t-1)+H+N\ge2nt-2$, hence $H+N\ge2n-2$.
At each genuine factor-fixing entry, average the D-AFT01
bounds over its active factors, and then enlarge the
denominator to $t$. This gives
$a^{t-1}(AI/t+B)$ with the entire positive constant $B$.
Scalars satisfy the same estimate at index zero.

At each full factor cycle the exact D-ATT08 estimate is
\[
 (1-1/t)e+\frac{a^{t-1}}t\nu(W_nx)-\frac{R_t(a)}t,
 \qquad R_t(a)=\sum_{j=1}^{t-1}a^{\gcd(t,j)}.            \tag{AFT.18}
\]
Only $\nu(W_nx)\ge AI(x)$ is used here, so identity cycle
products are allowed. The proved errors are $R_2(a)=a$
and $R_t(a)\le2a^{t-2}$ for $t\ge3$; in particular
$2R_t(a)/a^{t-1}\le2$ for all $t\ge2$.
There are at least ten factor-fixing entries, and summation
therefore yields
\[
 D-2e\ge\frac{a^{t-1}}t
       \bigl(2A(n-1)-2a+10tB-2\bigr)>B a^{t-1}>6.
                                                               \tag{AFT.19}
\]
For the strict comparison, the parenthesis is greater
than $tB$ because
\[
 9tB-\bigl(2a-2A(n-1)\bigr)-2
 \ge\left(\frac{36}{5}-\frac{101}{25}\right)S-2
 =\frac{79}{25}S-2>0.                                  \tag{AFT.20}
\]
This contradicts the exact Euler budget and completes
the all-count proof.
\end{proof}

\paragraph{Effective remaining-tail statement and strictness.}
For any fixed integer $K\ge3$, take the maximum of the
explicit $N_\mu$ over the finitely many partitions with
$3\le|\mu|\le K$. Above this effective maximum, the
non-Eisenstein transitive sources in D-AFT03 cannot have
tail size between three and $K$; D-AFT02 gives the same
statement for a single factor over every original
imaginary quadratic field. In the project's main
\emph{nonsplit} problem, combine this with the previously
proved standard and two-box exclusions in their stated
scopes. The tail size in any remaining non-Eisenstein
transitive ordinary alternating sequence with $n\to\infty$
must then tend to infinity. The single-factor standard
case still uses its separate actual-polarization splitness
gate. Transposing a label gives the same $A_n$
representation, so either long-row convention may be
used. No finite bound on $n$ is asserted when the tail
size itself grows.

For example, $\mu=(4)$ and any $n\ge N_{(4)}$ give a
label $(n-4,4)$ not among the previously excluded tails
of size at most three. Since $N_{(4)}\ge200$, its earlier
Gaussian long-row fraction
$4(n-5)/((n-1)(n-2))$ is below $1/10$; the earlier
necessary shape gate did not remove this entire family.
The new theorem removes it and every other fixed tail
eventually, without asserting that an actual source
previously existed. It supplies a new obstruction to
condition-A finite-cover suppliers, not a new algebraic
cycle or an unrefereed geometric premise.
